\section{\texorpdfstring{第~\ref{chap:differential-forms}~章}{微分形式}}\label{sec:solutions-differential-forms}
\subsection{\texorpdfstring{第~\ref{sec:tensors}~节}{张量}}\label{subsec:solutions-tensors}

\begin{description}[font=\normalfont,leftmargin=2em,labelwidth=1.5em]
\RGsolution{ex:alternating-projection-properties}{alternating-projection-properties}
设 $T\in\mathcal T^k(V^*)$.
\begin{enumerate}[label=(\alph*),ref=\alph*,leftmargin=*]
\item\label{sol:alternating-projection-image} 考虑张量 $\operatorname{Alt}(T):=\frac{1}{k!}\sum_{\sigma\in S_k}(\operatorname{sgn}\sigma)(T\circ\sigma)$. 若 $\sigma_0$ 是交换两个指标并固定其他指标的置换, 则
\begin{align*}
  \operatorname{Alt}(T)\circ\sigma_0
  &=\frac{1}{k!}\sum_{\sigma\in S_k}(\operatorname{sgn}\sigma)(T\circ\sigma\circ\sigma_0)\\
  &=\frac{1}{k!}(\operatorname{sgn}\sigma_0)
    \sum_{\sigma\in S_k}(\operatorname{sgn}(\sigma\circ\sigma_0))(T\circ(\sigma\circ\sigma_0))\\
  &=\frac{1}{k!}(\operatorname{sgn}\sigma_0)
    \sum_{\tau\in S_k}(\operatorname{sgn}\tau)(T\circ\tau)\\
  &=(\operatorname{sgn}\sigma_0)\operatorname{Alt}(T)=-\operatorname{Alt}(T),
\end{align*}
从而 $\operatorname{Alt}(T)$ 是交错张量.
\item\label{sol:alternating-projection-fixed-tensors} 若 $T$ 是交错的, 则对任意置换 $\sigma\in S_k$, 有 $T\circ\sigma=(\operatorname{sgn}\sigma)T$. 因此,
\[
  \operatorname{Alt}(T)=\frac{1}{k!}\sum_{\sigma\in S_k}(\operatorname{sgn}\sigma)(T\circ\sigma)
  =\frac{1}{k!}\sum_{\sigma\in S_k}(\operatorname{sgn}\sigma)^2T=T,
\]
这里使用了 $S_k$ 有 $k!$ 个元素.
\item 由 (\ref{sol:alternating-projection-image}) 知 $\operatorname{Alt}(T)$ 是交错张量, 所以由 (\ref{sol:alternating-projection-fixed-tensors}) 得 $\operatorname{Alt}(\operatorname{Alt}(T))=\operatorname{Alt}(T)$.
\end{enumerate}

\RGsolution{ex:prove-graded-commutativity}{graded-commutativity-of-wedge-product}
设 $\{T_1,\ldots,T_n\}$ 是 $V^*$ 的一组基. 有
\[
  T=\sum_{i_1<\cdots<i_k}a_{i_1\ldots i_k}T_{i_1}\wedge\cdots\wedge T_{i_k}
\]
及
\[
  S=\sum_{j_1<\cdots<j_m}b_{j_1\ldots j_m}T_{j_1}\wedge\cdots\wedge T_{j_m}.
\]
因此,
\begin{align*}
  T\wedge S
  &=\sum_{\substack{i_1<\cdots<i_k\\j_1<\cdots<j_m}}
     a_{i_1\ldots i_k}b_{j_1\ldots j_m}
     T_{i_1}\wedge\cdots\wedge T_{i_k}\wedge T_{j_1}\wedge\cdots\wedge T_{j_m}\\
  &=\sum_{\substack{i_1<\cdots<i_k\\j_1<\cdots<j_m}}
     a_{i_1\ldots i_k}b_{j_1\ldots j_m}(-1)^k
     T_{j_1}\wedge T_{i_1}\wedge\cdots\wedge T_{i_k}\wedge T_{j_2}\wedge\cdots\wedge T_{j_m}\\
  &=\sum_{\substack{i_1<\cdots<i_k\\j_1<\cdots<j_m}}
     a_{i_1\ldots i_k}b_{j_1\ldots j_m}(-1)^{km}
     T_{j_1}\wedge\cdots\wedge T_{j_m}\wedge T_{i_1}\wedge\cdots\wedge T_{i_k}\\
  &=(-1)^{km}S\wedge T.
\end{align*}

\RGsolution{ex:prove-pullback-of-tensors-properties}{pullback-of-tensors-properties}
\begin{enumerate}[label=(\arabic*),ref=\arabic*,leftmargin=*]
\item 若 $v_1,\ldots,v_{k+m}\in V$, 则
\begin{align*}
  &\bigl(F^*(T\otimes S)\bigr)(v_1,\ldots,v_{k+m})
  =(T\otimes S)(F(v_1),\ldots,F(v_{k+m}))\\
  &\quad=T(F(v_1),\ldots,F(v_k))S(F(v_{k+1}),\ldots,F(v_{k+m}))\\
  &\quad=(F^*T(v_1,\ldots,v_k))(F^*S(v_{k+1},\ldots,v_{k+m}))\\
  &\quad=((F^*T)\otimes(F^*S))(v_1,\ldots,v_{k+m}).
\end{align*}
\item 显然成立.
\item 若 $v_1,\ldots,v_{k+m}\in V$, 则
\begin{align*}
  &\bigl(F^*(T\wedge S)\bigr)(v_1,\ldots,v_{k+m})
  =(T\wedge S)(F(v_1),\ldots,F(v_{k+m}))\\
  &\quad=\frac{(k+m)!}{k!m!}(\operatorname{Alt}(T\otimes S))(F(v_1),\ldots,F(v_{k+m}))\\
  &\quad=\frac{1}{k!m!}\sum_{\sigma\in S_{k+m}}(\operatorname{sgn}\sigma)
       (T\otimes S)(F(v_{\sigma(1)}),\ldots,F(v_{\sigma(k+m)}))\\
  &\quad=\frac{1}{k!m!}\sum_{\sigma\in S_{k+m}}(\operatorname{sgn}\sigma)
       (F^*(T\otimes S))(v_{\sigma(1)},\ldots,v_{\sigma(k+m)})\\
  &\quad=\frac{1}{k!m!}\sum_{\sigma\in S_{k+m}}(\operatorname{sgn}\sigma)
       ((F^*T)\otimes(F^*S))(v_{\sigma(1)},\ldots,v_{\sigma(k+m)})\\
  &\quad=\frac{(k+m)!}{k!m!}(\operatorname{Alt}((F^*T)\otimes(F^*S)))(v_1,\ldots,v_{k+m})\\
  &\quad=((F^*T)\wedge(F^*S))(v_1,\ldots,v_{k+m}).
\end{align*}
\item 若 $v_1,\ldots,v_k\in V$, 则
\begin{align*}
  &\bigl((H\circ F)^*(T)\bigr)(v_1,\ldots,v_k)
  =T(H(F(v_1)),\ldots,H(F(v_k)))\\
  &\quad=(H^*T)(F(v_1),\ldots,F(v_k))=\bigl(F^*(H^*T)\bigr)(v_1,\ldots,v_k).
\end{align*}
\end{enumerate}
\end{description}
